\question
(0,5 балла) Вы проектируете сеть автобусных маршрутов в новом районе. Правительство даёт вам список остановок с указанием, сколько маршрутов должно проходить через конкретную остановку согласно нагрузке в районе: $(3, 3, 6, 6, 4, 4, 2)$. Вы сразу понимаете, что это вектор степеней, на основе которого можно построить граф.

Постройте граф согласно этому вектору, а также, чтобы ваша презентация выглядела убедительнее, постройте:

\begin{itemize}
\item частичный граф к исходному, который будет являться двудольным
\item подграф для построенного, такой чтобы при пересечении полученного подграфа и ранее построенного частичного графа получалось дерево;
\item надграф к исходному, такой, чтобы дополнительный граф к надграфу был регулярным;
\item 2 изоморфных графа с исходным, причем один с планарной укладкой, а другой с непланарной;
\item дополнительные графы к полученным выше частичному и подграфу относительно исходного графа;
\item нуль-граф к полученному выше надграфу.
\end{itemize}

Также отметьте на полученных графах мосты и точки сочленения.
